% JFQA Format Header - Converted from IEEE format
% Requirements: 12pt Times New Roman, double-spaced, 1" margins

\RequirePackage{silence}
\WarningFilter{latex}{Command \showhyphens has changed}

% ============================================
% JFQA-SPECIFIC PACKAGES
% ============================================

% Page layout: 8.5" x 11", 1" margins
\usepackage[margin=1in, letterpaper]{geometry}

% Double spacing
\usepackage{setspace}
\doublespacing

% Times New Roman font (JFQA requirement)
\usepackage{newtxtext,newtxmath}

% Author-year citations (JFQA style)
\usepackage[authoryear,round,sort&compress]{natbib}

% ============================================
% SECTION NUMBERING (JFQA Format)
% ============================================
% Sections: Roman numerals (I, II, III)
% Subsections: Letters (A, B, C)
% Subsubsections: Arabic numerals (1, 2, 3)

\renewcommand{\thesection}{\Roman{section}}
\renewcommand{\thesubsection}{\Alph{subsection}}
\renewcommand{\thesubsubsection}{\arabic{subsubsection}}

% Fix cross-references to use full JFQA format (e.g., "Section III.B.2")
\renewcommand{\thesubsection}{\thesection.\Alph{subsection}}
\renewcommand{\thesubsubsection}{\thesubsection.\arabic{subsubsection}}

% ============================================
% STANDARD PACKAGES (Compatible with JFQA)
% ============================================

\usepackage{amsmath,amssymb,amsfonts}
\usepackage{algorithmic}
\usepackage{graphicx}
\usepackage{textcomp}
\usepackage{xcolor}
\usepackage{booktabs}  % Professional tables
\usepackage{multirow}
\usepackage{url}
\usepackage[utf8]{inputenc}
\usepackage[T1]{fontenc}
\usepackage{gensymb}
\usepackage{microtype}  % Better typography
\usepackage[hypertexnames=false,plainpages=false]{hyperref}
\usepackage{array}
\usepackage{adjustbox}
\usepackage{float}
\usepackage{placeins}

% Table column spacing
\setlength{\tabcolsep}{5pt}

% ============================================
% JFQA EQUATION NUMBERING
% ============================================
% Equations numbered at left margin in parentheses
\makeatletter
\def\tagform@#1{\maketag@@@{(#1)\hfill}}
\makeatother
